\section{Additional Preliminaries}
\label{app:defs}

\subsection{Security assumptions for Lattices}\label{app:assumptions lattices}

Fix a public left factor \(F\), hint space \(H\), target space \(G\), predicate \(P\), and query bound \(Q\).
\begin{itemize}[leftmargin=1.25em]
  \item \textbf{\(\mathsf{Hint\mbox{-}vSIS}\).} The adversary chooses a query set \(\mathcal{Q}\subseteq H\) and a fresh
  target function \(g^\star\in G\setminus \mathcal{Q}\) before seeing the sampled matrix \(A\). It then receives short
  preimages of the queried hints \(q(A)\) under \(F(A)\) for every \(q\in\mathcal{Q}\), and must output a non-zero short
  vector \(u^\star\) such that \(F(A)u^\star=g^\star(A)\). The instance is restricted by the predicate condition
  \(P(F\cup \mathcal{Q}\cup(\{g^\star\}\setminus\{0\}))=1\).
  \item \textbf{\(s\)-\(\mathsf{Hint\mbox{-}vSIS}\).} This is the strong variant of the same experiment. The adversary still
  chooses the hint set \(\mathcal{Q}\) before seeing \(A\), but it is allowed to choose the fresh target \(g^\star\) only
  after seeing \(A\) and the hinted preimages.
  \item \textbf{\(s\)-\(\$\)-\(\mathsf{Hint\mbox{-}vSIS}\).} This is the strong random-hinted variant. It is the same as
  \(s\)-\(\mathsf{Hint\mbox{-}vSIS}\), except that the query set \(\mathcal{Q}\) is sampled uniformly at random from \(H\)
  rather than chosen by the adversary.
\end{itemize}

Now fix a keyed evaluation family \(f:K\times M\times X\to R_q^n\).
\begin{itemize}[leftmargin=1.25em]
  \item \textbf{\(\mathsf{GenISIS}_f\).} The challenger samples a public key \(\kappa\gets K\) and a uniform matrix
  \(A\getsunif R_q^{n\times m}\). It provides a preimage oracle that, on input \(\mu\in M\), samples \(\chi\gets X\),
  sets \(t=f(\kappa,\mu,\chi)\), returns a short preimage \(s\) such that \(As=t\), and records the tuple
  \((s,\mu,\chi)\). The adversary wins if it outputs a fresh tuple \((s^\star,\mu^\star,\chi^\star)\) not previously
  returned by the oracle such that \(As^\star=f(\kappa,\mu^\star,\chi^\star)\) and \(0<\|s^\star\|\le \beta\).
  \item \textbf{\(\mathsf{IntGenISIS}_f\).} This is an interactive version where the oracle includes the additional
  linear masking term and the stronger freshness condition used in DKLW's adaptive-security lifting.
\end{itemize}
The DKLW \(\mathsf{IntGenISIS}_f\)-to-\(\mathsf{GenISIS}_f\) reduction is used only under the parameter
conditions of their lifting theorem, including its ring, modulus, smoothing, dimension, and bound-loss
requirements~\cite[Theorem~6]{EPRINT:DKLW25}.


\subsection{Centering map and carries}
Let \(\BoundB\in\Z_{>0}\) and \(\Delta=2\BoundB+1\). Define the centering map
\(\centerfunc:[-2\BoundB,2\BoundB]\to[-\BoundB,\BoundB]\) by
\[
\centerfunc(x)=
\begin{cases}
x & \text{if } x\in[-\BoundB,\BoundB],\\
x+\Delta & \text{if } x\in[-2\BoundB,-\BoundB-1],\\
x-\Delta & \text{if } x\in[\BoundB+1,2\BoundB].
\end{cases}
\]
Applied coefficient-wise, this map allows unbiased combination of two bounded values when at least one is uniform.

\begin{lemma}[Centering preserves uniformity]\label{lem:center-uniform-app}
Fix any \(a\in[-\BoundB,\BoundB]\). If \(b\getsunif [-\BoundB,\BoundB]\) then \(\centerfunc(a+b)\) is uniform over
\([-\BoundB,\BoundB]\).
\end{lemma}

\subsection{Vanishing polynomials for membership constraints}
Membership constraints are expressed using vanishing polynomials over \(\Fq\) with signed lifting.

\begin{definition}[Vanishing polynomial for \(\lbrack -\BoundB,\BoundB\rbrack\)]\label{def:vanish-bound}
Let \(\BoundB\in\Z_{>0}\). Define
\[
  P_{\BoundB}(X) \;=\; \prod_{i=-\BoundB}^{\BoundB}\bigl(X-\langle i\rangle_q\bigr)\ \in\ \Fq[X].
\]
For \(x\in\Fq\), we have \(x\in[-\BoundB,\BoundB]\) (via signed lifting) iff \(P_{\BoundB}(x)=0\) in \(\Fq\).
\end{definition}

\begin{definition}[Vanishing polynomial for \(\{-1,0,1\}\)]\label{def:vanish-carry}
Define \(P_{1}(X)=(X+1)X(X-1)\in\Fq[X]\). Then \(P_1(x)=0\) iff \(x\in\{-1,0,1\}\).
\end{definition}
